\section{Ordering the Gaussian solid angles}
\label{vol:section}

This section gives the geometric step of the construction. All Gaussian
coordinates in this section have the same variance; by homogeneity we take
that variance to be one. The calculation concerns the probability of an
\emph{entire threshold function}, rather than the volume of a selected set of
its representations.

There are $k$ main bits $x$, $H\ge2$ helper bits $y$, and one selector bit $z$.
Write an affine score as
\[
 b+c z+\sum_{i=0}^{k-1}w_i x_i+\sum_{\ell=1}^H h_\ell y_\ell,
 \qquad u=b-c,\quad v=b+c.
\]
Thus $u$ and $v$ are independent Gaussian variables of variance two,
independent of the unit-variance coordinates $w_i,h_\ell$.
Let $M=2^k-1$, and write
\[
 R=\sum_iw_i,\qquad A_j=\sum_i\mathrm{bit}_i(j)w_i\quad(0\le j\le M).
\]
In particular $A_0=0$ and $A_M=R$.

Suppose a complete truth table on the branch $z=-1$ gives an open homogeneous
representation cone $\mathcal C$ in $(w,h,u)$. We use the following properties
of the calibration constructed in the preceding section. Every point of
$\mathcal C$ has $a=w_0>0$, satisfies
\begin{equation}\label{vol:binary}
 \left|\frac{A_j}{a}-j\right|<\epsilon\quad(0\le j\le M),
 \qquad h_\ell>R\quad(1\le\ell\le H),
\end{equation}
and can be parametrized by
\[
 h_\ell=H_0+\delta_\ell,\qquad \sum_\ell\delta_\ell=0,
 \qquad \theta=(w,\delta,u)\in\Omega_0,
 \qquad H_0>A(\theta).
\]
Here $\Omega_0$ is an open homogeneous cone, independent of $H_0$, and
$A$ is homogeneous of degree one. For a constant $B\ge1$,
\begin{equation}\label{vol:scaffold-bounds}
 R\le A(\theta)\le BR,\qquad
 F_0:=\sum_iw_i^2+\sum_\ell\delta_\ell^2+u^2/2\le B^2R^2.
\end{equation}
In particular, the condition on $H_0$ is exactly a lower ray, with no upper
restriction. All these statements concern every representation in the
calibration cone.

\subsection{The compatible functions are exactly the cutoff functions}

Let $q_j$ be the point on the branch $z=+1$ whose helper bits are all $+1$
and whose main bits encode the integer $j$. Add the labels $q_0\mapsto-1$
and $q_M\mapsto+1$ to the calibration. Assume $\epsilon<1/2$, so the values
$A_j$ increase strictly with $j$.
The negative label at $q_0$ gives $v+\sum h_\ell-R<0$.
If at least one helper bit is negative, the affine score is at most
\[
 v+\sum h_\ell-2\min_\ell h_\ell+R
 <2R-2\min_\ell h_\ell<0.
\]
Consequently all threshold functions compatible with the calibration are
exactly the $M$ functions $g_0,\ldots,g_{M-1}$: they agree with the calibrated
function on $z=-1$, are negative off the helper face on $z=+1$, and satisfy
\[
 g_j(q_i)=-1\ \ (i\le j),\qquad
 g_j(q_i)=+1\ \ (i>j).
\]
Their full representation cones are
\begin{equation}\label{vol:full-cell}
 (w,h,u)\in\mathcal C,\qquad
 -\sum h_\ell+R-2A_{j+1}<v<-\sum h_\ell+R-2A_j.
\end{equation}
Each cell has positive Gaussian measure: $\mathcal C$ is nonempty and open,
and every displayed interval has positive length. Denote its Gaussian
probability by $p_j$.

\subsection{Radial integration and the common helper direction}

Set $D=k+H+1$ and $s=(D+1)/2$. Normalize $a=w_0$ to one.
There are $D+1$ original affine coefficients, and the radial change of
variables has Jacobian $a^D$. Integrating $a>0$ in the Gaussian density gives
a common positive constant times the power $-s$ of the normalized Gaussian
quadratic form. The changes $(b,c)\mapsto(u,v)$ and
$h\mapsto(H_0,\delta)$ contribute only constant Jacobians.
For fixed normalized $\theta$ put
\begin{equation}\label{vol:psi}
 \psi_\theta(r)=\int_{A(\theta)}^\infty
 \left[F_0+HH_0^2+\frac{(HH_0+r)^2}{2}\right]^{-s}\,dH_0.
\end{equation}
Equation~\eqref{vol:full-cell}, with $-v=HH_0+r$, therefore gives
\begin{equation}\label{vol:mass-integral}
 p_j=K\int_{\{\theta\in\Omega_0:w_0=1\}}
 \int_{-R+2A_j}^{-R+2A_{j+1}}\psi_\theta(r)\,dr\,d\theta,
\end{equation}
where $K>0$ and the cross-section measure are the same for every $j$.
All inner integration intervals lie in $[-R,R]$.

Define constants
\begin{align*}
 C_1&=B^2+H(B+1)^2+\frac{[H(B+1)+1]^2}{2},\\
 T_0&=2\max\{1,C_1/H\},\\
 C_2&=B^2+HT_0^2+\frac{(HT_0+1)^2}{2},\\
 \ell&=\frac{s(H-1)}{2C_2(M+1)}.
\end{align*}
The constants are polynomial in $H,B$, apart from the explicit factor
$(M+1)^{-1}$ in $\ell$.

\begin{lemma}[A uniform slope after integrating the common direction]
\label{vol:slope}
For every normalized $\theta$ and every $r\in[-R,R]$, if $\epsilon\le1$,
\[
 -\frac{\psi_\theta'(r)}{\psi_\theta(r)}\ge\ell.
\]
\end{lemma}
\begin{proof}
Write $F=F_0+HH_0^2+(HH_0+r)^2/2$.
The part of the integral in~\eqref{vol:psi} with
$A\le H_0\le A+R$ is at least $R(C_1R^2)^{-s}$.
The part with $H_0>T_0R$ is at most
\[
 \int_{T_0R}^\infty(HH_0^2)^{-s}\,dH_0
 =\frac{H^{-s}(T_0R)^{1-2s}}{2s-1}.
\]
The ratio of this upper bound to the preceding lower bound is
\[
 \frac{(C_1/H)^sT_0^{1-2s}}{2s-1}\le\frac12.
\]
Indeed $s\ge1$, and substituting $T_0=2\max\{1,C_1/H\}$ proves the inequality
separately for $C_1/H\le1$ and $C_1/H\ge1$.
Thus at least half the probability measure with density proportional to
$F^{-s}$ lies in $H_0\le T_0R$.
Differentiation under the integral is justified by the integrable power
bounds at infinity and gives
\[
 -\frac{\psi_\theta'(r)}{\psi_\theta(r)}
 =s\,\mathbb E\left[\frac{HH_0+r}{F}\right].
\]
The integrand is positive everywhere, because $H_0>A\ge R$ and $H\ge2$.
On $H_0\le T_0R$, its numerator is at least $(H-1)R$ and its denominator
at most $C_2R^2$. Hence the last display is at least
$s(H-1)/(2C_2R)$. Finally~\eqref{vol:binary} gives
$R=A_M<M+\epsilon\le M+1$.
\end{proof}

\begin{proposition}[Quantitative ordering of whole function cells]
\label{vol:ratio}
If
\[
 \epsilon\le\frac18,\qquad \epsilon\le\frac{\ell}{32},
\]
then, for $0\le j<M-1$,
\[
 p_j\ge(1+\ell/4)p_{j+1}.
\]
\end{proposition}
\begin{proof}
For fixed $\theta$, let $l_j=2(A_{j+1}-A_j)$.
By~\eqref{vol:binary}, $2-4\epsilon<l_j<2+4\epsilon$;
in particular $3/2\le l_j\le5/2$ and $l_{j+1}-l_j\le8\epsilon$.
Let $t$ be the common boundary of the two adjacent integration intervals.
Lemma~\ref{vol:slope} implies
\[
 \int_{t-l_j}^{t}\psi_\theta(r)\,dr
 \ge\psi_\theta(t)(l_j+\ell l_j^2/2),
 \qquad
 \int_t^{t+l_{j+1}}\psi_\theta(r)\,dr
 \le\psi_\theta(t)l_{j+1}.
\]
The difference of the first mass and the second is at least
\[
 \psi_\theta(t)(\ell l_j^2/2-8\epsilon)
 \ge(7\ell/8)\psi_\theta(t),
\]
whereas the second mass is at most $(5/2)\psi_\theta(t)$.
Thus the first is at least $(1+7\ell/20)$ times the second, which is stronger
than the asserted factor $1+\ell/4$. Integrate this inequality
in~\eqref{vol:mass-integral}.
\end{proof}

\subsection{A quantitative width bound in the cold posterior}

Take the target to be $g_{M-1}$ and, after calibration, reveal
$q_1,q_2,\ldots,q_{M-1}$, each with label $-1$.
There are $T=M-1=2^k-2$ designated rounds.
Suppose the cold posterior on surviving threshold functions is proportional
to $p_j^N$, as established by the posterior calculation for sign networks.
Immediately before $q_i$ is revealed, the survivors are
$g_{i-1},\ldots,g_{M-1}$. Exactly $g_{i-1}$ predicts the wrong label $+1$.
Writing $\rho=1+\ell/4$, Proposition~\ref{vol:ratio} and a geometric-series
bound give
\begin{equation}\label{vol:prediction-bound}
 \Pp(\text{correct label at }q_i\mid\text{past labels})\le\rho^{-N}.
\end{equation}
For completeness, if the surviving weights are $b_0,b_1,\ldots$ with
$b_{r+1}\le q b_r$ and $q=\rho^{-N}$, their tail has mass at most
$qb_0/(1-q)$, so its posterior fraction is at most $q$.
Consequently
\[
 N\ge\frac{\log(T/\eta)}{\log(1+\ell/4)}
\]
ensures, by a union bound, that sampled posterior predictions make every one
of these $T$ mistakes with probability at least $1-\eta$.
The same estimate implies a wrong posterior majority at every designated
round whenever $\rho^{-N}<1/2$.
When $H$ and $B$ are polynomial in $k$, the displayed cold width threshold is
$2^k\operatorname{poly}(k)\log(T/\eta)$.
This quantitative threshold concerns the cold posterior. A simultaneous
width-and-scale limit also needs the separate finite-scale posterior error
bound; its dimension-dependent constants should not be suppressed into this
cold width estimate.
